% Chapter 8, generated by make_chapters.py from papers/Ch08_Bio_12_Frontier_Elimination/anccft25.tex.
% Source paper: Bio 12
\chapter[Frontier elimination]{The double-stranded lattice sum by frontier elimination}\label{chap:ch08}
\begin{chaptersummary}
Chapter~\ref{chap:ch07} wrote the number of double-stranded reconstructions of a circular sequence as a sum over
the lattice points of a box, $\#\DS(S,k)=\sum_f\Nseq\big((m+f)/2\big)$ with $f$ running over
$\Lm\cap\{|f|\le m,\ f\equiv m\ (2)\}$, and enumerated that box by backtracking. The
enumeration works at $\rank\Lm=4$ and $12$ and dies at $48$: $\phi$X174 at $k=10$ is the blank
row of Table~\ref{ch07:tab:lattice}. The obstruction is not the rank. The circulation condition is one
linear equation per vertex, and each equation involves only the four endpoints of one
$\rho$-orbit, so the orbits can be processed in an order while carrying, as state, the partial
imbalance at the vertices still open --- a product of sparse transfer operators whose cost is
governed by the width of the \emph{frontier}, not by the dimension of the lattice. At $k=10$,
where the box lies in a $48$-dimensional lattice and has more than two billion points, the
elimination needs $50\,758$ states and $8.4$ seconds: $\phi$X174 has exactly
$2\,481\,687\,900$ box points at $k=10$. Carrying in addition the partition of the frontier
into the classes already joined counts the points with connected support, of which there are
$909\,103\,210$; both columns Chapter~\ref{chap:ch07} left empty are now filled, and only $\#\DS$ itself is not
(Chapter~\ref{chap:ch09} fills it).
Two further things follow. Connectivity of $\supp c$,
carried in Chapter~\ref{chap:ch07} as a separate filter, is not a side condition at all: a balanced digraph has
strongly connected weak components, so its out-Laplacian has nullity equal to the number of
them, and $t(G_c)$ --- hence $\Nseq(c)$ --- vanishes on a disconnected support of its own
accord. And the order that works is not the genome's own: along the reference coordinate the
frontier at $k=10$ is $62$ wide and the elimination does not finish, while a greedy order is
$26$ wide, which is the interleaving of repeats of Theorem~\ref{ch01:thm:interleave} seen as a width.
\end{chaptersummary}

\section{What is being counted}\label{ch08:sec:setup}

We keep the notation of Chapter~\ref{chap:ch07} exactly: $D_k=D_k(S)$ is the labelled double
cover of Definition~\ref{ch07:def:cover}, $A$ its set of contents, $m(a)$ the number of labelled
arcs of content $a$, $\rho$ the reverse-complement involution, and $\Lm$ the $(-1)$-eigenlattice
of $\rho$ on the circulations of $D_k$. A double-stranded reconstruction is as in
Definition~\ref{ch07:def:ds}, and by Theorem~\ref{ch07:thm:lattice}
\begin{equation}\label{ch08:eq:sum}
\#\DS(S,k)\;=\;\sum_{f}\Nseq\Big(\frac{m+f}{2}\Big),
\qquad
\Nseq(c)=\frac{t(G_c)\,\Loc(G_c)}{\prod_a c(a)!},
\end{equation}
the sum over $f\in\Lm$ in the box $|f|\le m$, $f\equiv m\ (2)$, with $G_c$, $t(G_c)$ and
$\Loc(G_c)$ as there; the division by $\prod_a c(a)!$ passes from labelled arcs to sequences
(Remark~\ref{ch07:rem:norm}).

Chapter~\ref{chap:ch07} enumerated the box by backtracking over a basis of $\Lm$ with capacity pruning, and
recorded the cost: $10$ points at $k=12$ where $\rank\Lm=4$, $528$ at $k=11$ where
$\rank\Lm=12$, and nothing at $k=10$, where $\rank\Lm=48$ and the search does not return.

\section{The constraint system and its frontier}\label{ch08:sec:frontier}

The point of this chapter is that \eqref{ch08:eq:sum} is indexed by the solutions of a very sparse
linear system, and that sparsity, not the rank, is what governs the count.

\begin{definition}[the orbit variables]\label{ch08:def:vars}
$\rho$ acts on $A$ without fixed points except for the self-reverse-complementary contents,
which carry $f(a)=0$ because $\rho f=-f$. Choose one representative $a$ from each free orbit
$\{a,\rho a\}$ and let $g_j=f(a_j)$, so $f(\rho a_j)=-g_j$. The box condition is
$g_j\in\{-m_j,-m_j+2,\dots,m_j\}$ with $m_j=m(a_j)$, and $f$ is a circulation exactly when,
for every vertex $v$,
\begin{equation}\label{ch08:eq:bal}
\sum_j B(v,j)\,g_j=0,\qquad
B(v,j)=[v=\mathrm{t}(a_j)]-[v=\mathrm{h}(a_j)]-[v=\mathrm{t}(\rho a_j)]+[v=\mathrm{h}(\rho a_j)] .
\end{equation}
\end{definition}

The essential feature is that column $j$ of $B$ is supported on at most four vertices --- the
two endpoints of $a_j$ and the two of $\rho a_j$ --- whatever the rank of $B$ may be.

\begin{definition}[frontier]\label{ch08:def:front}
Fix an order $\pi=(j_1,\dots,j_R)$ of the orbits. A vertex $v$ is \emph{open after step $i$}
if some $j_\ell$ with $\ell\le i$ and some $j_\ell$ with $\ell>i$ both have $B(v,j_\ell)\ne0$.
Write $\Fr_i(\pi)$ for the set of open vertices and $w(\pi)=\max_i|\Fr_i(\pi)|$ for the
\emph{width} of $\pi$. The minimum of $w(\pi)$ over all orders is the pathwidth of the
hypergraph whose vertices are the vertices of $D_k$ and whose hyperedges are the supports of
the columns of $B$.
\end{definition}

\begin{theorem}[frontier elimination]\label{ch08:thm:elim}
Let $\kappa_i(v)=\sum_{\ell>i}|B(v,j_\ell)|\,m_{j_\ell}$ be the capacity still available at
$v$ after step $i$, and let
$X_i=\prod_{v\in\Fr_i(\pi)}\{-\kappa_i(v),\dots,\kappa_i(v)\}$.
Define $T_i\colon \Z^{X_{i-1}}\to\Z^{X_i}$ by
\[
(T_i\phi)(y)\;=\;\sum_{\substack{g\in\{-m_{j_i},\dots,m_{j_i}\}\\ g\equiv m_{j_i}(2)}}
\ \sum_{\substack{x\in X_{i-1}\\ x+gB(\cdot,j_i)\ \text{restricts to}\ y\\
\text{and vanishes at every vertex closing at }i}}\phi(x).
\]
Then
\[
\big|\Lm\cap\{|f|\le m,\ f\equiv m\ (2)\}\big|\;=\;(T_R\cdots T_1\delta_0)(\,)\;,
\]
and the whole product costs $O\!\left(\sum_i |X_{i-1}|\,(m_{j_i}+1)\,|\Fr_i|\right)$ integer
operations, with
\[
|X_i|\;\le\;\prod_{v\in\Fr_i}\big(2\kappa_i(v)+1\big)
\;\le\;\big(2\max_v\kappa_i(v)+1\big)^{w(\pi)} .
\]
\end{theorem}

\begin{proof}
A point of the box is an assignment $g\in\prod_j\{-m_j,\dots,m_j\}$ of the right parity
satisfying \eqref{ch08:eq:bal}. Process the orbits in the order $\pi$ and record, after step $i$,
the partial sums $x(v)=\sum_{\ell\le i}B(v,j_\ell)g_{j_\ell}$ at the open vertices. A vertex
$v$ that is not open after step $i$ receives no further contribution, so \eqref{ch08:eq:bal} at $v$
is decided then and there: it is the condition, imposed in $T_i$, that $x(v)=0$ at every
closing vertex. A vertex not yet touched has $x(v)=0$ and need not be recorded. Two
assignments of $g_{j_1},\dots,g_{j_i}$ with the same restriction to $\Fr_i$ therefore admit
exactly the same completions, which is the Markov property that makes the transfer operators
well defined and their product the count. The capacity bound holds because
$|x(v)|=\big|\sum_{\ell>i}B(v,j_\ell)g_{j_\ell}\big|\le\kappa_i(v)$ for any state that has a
completion, so states outside $X_i$ may be dropped.
\end{proof}

\begin{remark}[what the cost does and does not depend on]\label{ch08:rem:cost}
Theorem~\ref{ch08:thm:elim} contains no reference to $\rank\Lm$. The backtracking of Chapter~\ref{chap:ch07}
has cost proportional to the number of \emph{nodes} of a search tree of depth $R$, which for a
lattice of rank $r$ is at least the number of points and so at least exponential in $r$; the
elimination has cost proportional to the number of \emph{states}, and a state is a vector of
imbalances on the frontier. At $k=10$ these are $2\,481\,687\,900$ and $50\,758$.
\end{remark}

\begin{proposition}[product-form weights]\label{ch08:prop:weights}
Let $u_j$ be functions on the values of $g_j$ and $h_v$ functions on $\Z_{\ge0}$, and
\[
w(c)=\prod_j u_j(g_j)\prod_v h_v\big(d^+_c(v)\big).
\]
Then $\sum_f w(c_f)$ over the box is computed by Theorem~\ref{ch08:thm:elim} after enlarging the
state at each open vertex by its accumulated out-degree. In particular
$\Loc(G_c)$ and $\prod_a c(a)!$ are of this form.
\end{proposition}

\begin{proof}
$d^+_c(v)=\sum_{a:\,\mathrm{t}(a)=v}c(a)$ accumulates one orbit at a time exactly as $x(v)$
does, and is final when $v$ closes, where $h_v$ may be applied; $u_j$ multiplies into the
transfer operator at step $j$. $\Loc$ takes $h_v(d)=(d-1)!$ and $u_j=1$, while
$\prod_a c(a)!$ takes $h_v=1$ and $u_j(g)=\big(\tfrac{m_j+g}{2}\big)!\big(\tfrac{m_j-g}{2}\big)!$.
\end{proof}

\section{Connectivity is the matrix-tree determinant}\label{ch08:sec:conn}

Chapter~\ref{chap:ch07} carried, beside the box conditions, a predicate: the support of $c$ must be connected,
or the point contributes nothing. The predicate is real but it is not an extra condition.

\begin{lemma}[connectivity is not a side condition]\label{ch08:lem:conn}
Let $c\in\Z_{\ge0}^A$ be balanced, $V_c=\supp c$ the set of vertices it meets, $L_c$ the
out-degree Laplacian of $G_c$ on $V_c$ and $\widetilde L_c$ the matrix $L_c$ with one row and
the corresponding column deleted. Then $\det\widetilde L_c=t(G_c)$, and
\[
\det\widetilde L_c=0\iff G_c\ \text{is disconnected.}
\]
Consequently $\Nseq(c)=0$ whenever $\supp c$ is disconnected, and the connectivity predicate
may be deleted from \eqref{ch08:eq:sum} without changing the sum.
\end{lemma}

\begin{proof}
$\det\widetilde L_c=t(G_c)$ is the matrix-tree theorem for digraphs \cite{Tutte48}. A balanced
digraph has $d^+=d^-$ at every vertex, so each of its weakly connected components is Eulerian
and hence strongly connected; therefore the kernel of $L_c$ is spanned by the indicators of
the components and $\dim\ker L_c=r$, the number of them. Deleting one row and one column drops
the nullity by at most one, so $r\ge2$ forces $\det\widetilde L_c=0$, while $r=1$ gives
$t(G_c)\ge1$.
\end{proof}

\begin{remark}\label{ch08:rem:predicate}
The gain is conceptual rather than computational: as a filter, connectivity is a union-find
pass and the determinant is a cubic elimination, so one still tests connectivity first. What
Lemma~\ref{ch08:lem:conn} removes is the impression that \eqref{ch08:eq:sum} is a lattice sum
\emph{subject to a combinatorial constraint}. It is a lattice sum, full stop; the constraint
is already inside $\Nseq$. This is also the answer, in the only form it has, to the natural question of an
algebraic criterion of connected support avoiding graph traversal: the
criterion is the vanishing of a determinant whose entries are linear in $c$.
\end{remark}

Connectivity can also be carried \emph{inside} the elimination, and then it counts.

\begin{theorem}[connected points]\label{ch08:thm:connected}
Enlarge the state of Theorem~\ref{ch08:thm:elim} by (i) the partition of $\Fr_i$ into the classes
connected by arcs already chosen, (ii) for each class, whether it meets a chosen arc, and
(iii) the number of classes that have left the frontier while carrying a chosen arc. Prune
whenever (iii) exceeds one, and accept only states with (iii) equal to one and empty frontier.
The resulting count is the number of box points with connected support. The state count is at
most $|X_i|\cdot B(|\Fr_i|)\cdot2^{|\Fr_i|}\cdot2$ with $B$ the Bell numbers.
\end{theorem}

\begin{proof}
A class leaves the frontier only when no vertex of it is open, and after that no arc can join
it to anything; so a class that leaves while meeting a chosen arc is a finished connected
component of $\supp c$, and conversely every component is finished exactly once, at the step
where its last vertex closes. A support is connected iff it has exactly one component, which
is condition (iii); classes meeting no chosen arc are not vertices of $\supp c$ and must be
ignored, which is what (ii) records. Since the class of a vertex and its flag are determined
by the chosen arcs among the processed orbits, the state is again Markov.
\end{proof}

\begin{remark}[the connectivity state is cheaper than it looks]\label{ch08:rem:bell}
The Bell factor in Theorem~\ref{ch08:thm:connected} is a gross overestimate, because the rule
``at most one finished component'' forbids most partitions from ever arising: at $k=12$ and
$k=11$ the connected elimination peaks at $8$ and $136$ states against $4$ and $20$ for the
plain one. It is nevertheless the factor that decides the cost at $k=10$, where a frontier of
$26$ vertices is small for imbalances and large for partitions: $9\,749\,212$ states against
$50\,758$, and seventy-eight minutes against eight seconds. $B(26)$ is $5\times10^{19}$, so
what is actually reached is a vanishing fraction of the bound.
\end{remark}

\section{The genome's own coordinate is not the good order}\label{ch08:sec:width}

Theorem~\ref{ch08:thm:elim} turns the problem into the choice of $\pi$. There is an obvious
candidate, and it is the one the biology suggests: take the orbits in the order of their
position along the reference genome, so that the elimination sweeps the chromosome once. It
is bounded, and it is not good.

\begin{proposition}[the coordinate order]\label{ch08:prop:coord}
Order the orbits by the position of their first occurrence on $S$. Then $\Fr_i$ is contained
in the set of repeated $k$-mers having occurrences in orbits on both sides of step $i$; in
particular, if no more than $\omega$ repeat classes straddle any coordinate, the elimination
runs in $O\big(R\,(2\max\kappa+1)^{\omega}\big)$ and $\#\DS$'s box is counted in time
polynomial in $N$ for fixed $\omega$.
\end{proposition}

\begin{proof}
Immediate from Definition~\ref{ch08:def:front}: the vertices met by orbit $j$ are endpoints of
$a_j$ and $\rho a_j$, and a vertex is open only between the first and last orbit meeting it.
\end{proof}

The hypothesis is the point. Table~\ref{ch08:tab:width} gives the widths.

\begin{table}[t]
\centering
\small
\begin{tabular}{rrrrrrr}
\toprule
$k$ & $|V|$ & contents & orbits & $\rank\Lm$ & coordinate width & greedy width\\
\midrule
12 &   7 &  14 &   7 &  4 &  6 &  6\\
11 &  30 &  53 &  26 & 12 & 18 & 10\\
10 & 125 & 222 & 111 & 48 & 62 & 26\\
\bottomrule
\end{tabular}
\caption{The width of the two orders. Along the genome the frontier at $k=10$ is $62$ vertices
wide and the elimination does not finish --- it passes $1.9$ million states and is abandoned;
the greedy order is $26$ wide and finishes in $8.4$ seconds with $50\,758$ states. Check (O).}
\label{ch08:tab:width}
\end{table}

\begin{remark}[why the coordinate order fails]\label{ch08:rem:interleave}
A vertex stays open from its first to its last incident orbit, so the coordinate width at a
position $x$ counts the repeat classes straddling $x$. Two things inflate it. First,
$\rho$ pairs an arc of $S$ with an arc of $\bar S$, and the two lie at reflected coordinates,
so an orbit is not local on the reference at all: every orbit near the origin of $S$ is also an
orbit near the end of $\bar S$. Second, and this is the substance, inverted repeats of a real
sequence \emph{interleave} --- that is the content of the interleaving theorem (Theorem~\ref{ch01:thm:interleave}),
where the crossing pattern of repeats is exactly what the sandpile group of the de Bruijn graph
records. Interleaved repeats are nested intervals crossing each other, and a family of
pairwise crossing intervals forces every one of them to be open at the crossing point. The
coordinate width is therefore a measure of interleaving, and a genome with interleaved
inverted repeats has no narrow sweep. The greedy order abandons the chromosome and follows the
constraint hypergraph instead.
\end{remark}

\section{\texorpdfstring{$\phi$X174 at $k=10$}{phiX174 at k=10}}\label{ch08:sec:phix}

\begin{table}[t]
\centering
\small
\begin{tabular}{rrrrrrrr}
\toprule
& & \multicolumn{2}{c}{box} & \multicolumn{2}{c}{connected} & &\\
\cmidrule(lr){3-4}\cmidrule(lr){5-6}
$k$ & $\rank\Lm$ & points & states & points & states & \%\ conn. & $\#\DS$\\
\midrule
13 & 0 & 1 & --- & 1 & --- & 100 & 1\\
12 & 4 & 10 & 4 & 8 & 8 & 80.0 & 10\\
11 & 12 & 528 & 20 & 218 & 136 & 41.3 & 864\\
10 & 48 & \textbf{2\,481\,687\,900} & 50\,758 & \textbf{909\,103\,210} & 9\,749\,212 & 36.6 & ---\\
\bottomrule
\end{tabular}
\caption{The blank row of Table~\ref{ch07:tab:lattice}, filled in both of its counting columns. The
rows for $k=12$ and $k=11$ reproduce that table exactly, by the elimination and independently
by the backtracking enumerator; the box count at $k=10$ is confirmed by pinning one coordinate
and summing the elimination over its values. ``states'' is the peak number of frontier states:
$8.4$ seconds for the box, $78$ minutes and $10$\,GB for the connected count.
Checks (T), (K), (W), (M).}
\label{ch08:tab:phix}
\end{table}

$\phi$X174 at $k=10$ has $125$ vertices, $222$ contents, $111$ free $\rho$-orbits and
$\rank\Lm=48$; the box
$\Lm\cap\{|f|\le m,\ f\equiv m\ (2)\}$ contains exactly
\[
2\,481\,687\,900
\]
points. Three remarks on the number. It is exact and integral throughout --- the elimination
adds integers and never divides. It is confirmed independently: pinning one orbit variable to
each of its admissible values in turn and summing the resulting eliminations, for two different
orbits, returns the same total, and the two orbits have disjoint schedules. And it is not
reachable by enumeration: $2.48\times10^9$ points at the rate at which Chapter~\ref{chap:ch07} enumerates
$528$ of them would be several months, and the relevant comparison is that the whole count
took $8.4$ seconds because it never listed a point.

\begin{remark}[the connected fraction]\label{ch08:rem:frac}
The proportion of box points with connected support falls with $k$: $100\%$, $80\%$, $41.3\%$,
$36.6\%$ at $k=13,12,11,10$. This is what one expects and it is worth saying why. A point
chooses one arc from each free $\rho$-orbit, so the support of $c$ has close to $|A|/2$ arcs on
$|V|$ vertices whatever the point; at $k=10$ that is about $111$ arcs on $125$ vertices, which
is sparse enough that a random choice usually falls apart. What the elimination shows is that
$36.6\%$ of them do not, and that the surviving set is still of size $9\times10^8$ --- so the
failure to reach $\#\DS$ at $k=10$ is not for want of a small enough index set.
\end{remark}

\begin{remark}[what is still missing, exactly]\label{ch08:rem:wall}
Of the three counting columns of Table~\ref{ch08:tab:phix} only $\#\DS$ is now empty at $k=10$, and
it is empty for a reason of a different kind from the other two. Those were questions of state
size: Theorem~\ref{ch08:thm:elim} settles the first in seconds and Theorem~\ref{ch08:thm:connected} the
second in an hour and a quarter.

The $\#\DS$ column is a question of principle. By
Proposition~\ref{ch08:prop:weights} the elimination sums any weight that factors into per-orbit and
per-vertex terms, and two of the three factors of $\Nseq$ do: $\Loc(G_c)$ is a product over
vertices of a function of the out-degree, and $\prod_a c(a)!$ a product over orbits. The third,
$t(G_c)$, does not factor over any bounded frontier --- it is a determinant of order up to
$|V_c|$, a global function of the whole support --- and it is the one that carries the count.
Evaluating it point by point is cheap per point and hopeless in aggregate: $9.09\times10^8$
determinants of order up to $125$. So $\#\DS(\phi X174,10)$ remains out of reach of the elimination (Chapter~\ref{chap:ch09} obtains it
point by point), and it is
now clear that what stands in the way is the arborescence factor and not the lattice --- the
index set is known exactly, and it is the summand that cannot be carried. Table~\ref{ch07:tab:lattice}
gives the two ends of the interval it lies in: the single-strand count
$\Nseq(S)+\Nseq(\bar S)=8\,610\,708\,632$ is a lower bound, since $\spec(S)$ and $\spec(\bar S)$
are two of the connected points, and $\mathrm{ec}(D_{10})=1.368\times10^{36}$ is an upper one.
\end{remark}

\section{Verification}\label{ch08:sec:checks}

\texttt{ds\_transfer\_matrix.py} is standalone --- Python and the standard library, no
\texttt{numpy}, no \texttt{sympy} and no other script of the book --- and rebuilds $D_k$ from
the definition rather than importing it, so that check (D) is a genuine cross-check of Table~\ref{ch07:tab:rank} and not a tautology. Exact integer arithmetic throughout; the determinants of
Lemma~\ref{ch08:lem:conn} are computed over $\mathbb Q$ with \texttt{fractions}.

\begin{table}[t]
\centering
\small
\begin{tabular}{@{}l@{\ \ }p{0.88\textwidth}@{}}
\toprule
key & claim\\
\midrule
(D) & the double cover rebuilt from the definition reproduces $|V|$, $\#\mathrm{fix}_V$,
$|A|$, $\#\mathrm{fix}_A$, $\sum m$ and $\rank\Lm$ of Table~\ref{ch07:tab:rank} at $k=10,11,12$\\
(T) & the frontier elimination and the backtracking enumerator return the same box count,
$10$ and $528$\\
(C) & on all $538$ box points at $k=12,11$: $\supp c$ connected $\iff$
$\det\widetilde L_c\ne0$ --- Lemma~\ref{ch08:lem:conn}\\
(W) & summing $\Nseq$ over the box returns $8$ and $218$ connected points and
$\#\DS=10$ and $864$\\
(O) & coordinate width $6,18,62$ against greedy width $6,10,26$; at $k=10$ the coordinate
order does not finish\\
(K) & $\phi$X174 at $k=10$: $2\,481\,687\,900$ box points, and the same total after pinning
each of two orbits\\
(N) & the connectivity-class elimination returns Chapter~\ref{chap:ch07}'s $8$ and $218$\\
(M) & $\phi$X174 at $k=10$: $909\,103\,210$ box points with connected support\\
\bottomrule
\end{tabular}
\caption{The eight checks. All pass. Check (M) takes about eighty minutes and $10$\,GB; the
flag \texttt{-{}-quick} skips it and the rest of the battery runs in six minutes.}
\label{ch08:tab:checks}
\end{table}

\section{Scope, provenance, and what is open}\label{ch08:sec:scope}

\paragraph{Scope.} Theorem~\ref{ch08:thm:elim}, Proposition~\ref{ch08:prop:weights},
Lemma~\ref{ch08:lem:conn} and Theorem~\ref{ch08:thm:connected} hold for any circular $S$ and any $k$.
Proposition~\ref{ch08:prop:coord} is a bound, and Table~\ref{ch08:tab:width} shows it is not a good one
on real sequence. Section~\ref{ch08:sec:phix} is a computation on one named genome and claims
nothing beyond it. Nothing here improves \eqref{ch08:eq:sum} itself, which is Chapter~\ref{chap:ch07}'s.

\paragraph{Provenance.} The lattice sum, $\Lm$, its rank by the Hopf trace formula, and the
$k=12$ and $k=11$ values are Chapter~\ref{chap:ch07}; the double cover and its balance are Chapter~\ref{chap:ch05},
its deck involution as an anti-automorphism Theorem~\ref{ch06:thm:anti}, the interleaving of
repeats Theorem~\ref{ch01:thm:interleave}. The BEST theorem is \cite{BEST,SmithTutte} and the matrix-tree theorem for
digraphs is \cite{Tutte48}. Frontier elimination is not new: it is the standard
dynamic-programming-over-a-path-decomposition, and in the form used in
Theorem~\ref{ch08:thm:connected} --- carrying a partition of the frontier so as to count connected
subgraphs --- it is the frontier-based search of \cite{KIIM17}, itself descended from Knuth's
\textsc{Simpath} \cite{Knuth4A}; the width parameter is pathwidth \cite{Bodlaender98}. What is
specific here is the observation that the double-stranded reconstruction count is indexed by a
system whose columns have support four, so that the frontier is narrow for reasons having
nothing to do with $\rank\Lm$, and the identification of the coordinate order as the wrong one.
For counting lattice points in a polytope of \emph{fixed} dimension there is Barvinok's
algorithm \cite{Barvinok94}; it does not apply here, the dimension being the thing that grows.

\paragraph{Open.} \emph{First}, $\#\DS(\phi X174,10)$, for which something must be done about
$t(G_c)$ (Chapter~\ref{chap:ch09} computes it by a different route); the natural attempt is to carry a partial Schur complement of $L_c$ on the frontier,
which is a matrix state and not a scalar one, and whether that can be made to fit is the real
question this chapter leaves. An order chosen to keep partitions few rather than imbalances small
would help both it and Theorem~\ref{ch08:thm:connected}, whose cost at $k=10$ is entirely the
partition factor. \emph{Second}, the width of a bacterial genome. The
constraint hypergraph of \emph{E.~coli} at $k=31$ falls apart into many small blocks, one per
repeat cluster, and the elimination factors over them; whether the largest block stays narrow
at the word lengths where the double cover is interesting is not known to us. \emph{Third},
whether the coordinate width of Proposition~\ref{ch08:prop:coord} is, for a random sequence with a
prescribed inverted-repeat structure, a computable function of that structure --- it would
make ``how hard is double-stranded assembly here'' a statistic of the genome.
